\section{Logical Implication, Necessary Conditions, and Sufficient Conditions}

\subsection{Introduction}

Logic is the language we use to organize mathematical reasoning. Before proving
theorems, we need to understand how mathematical statements are connected.

A \emph{statement} is a sentence that is either true or false, but not both.

For example,

\[
\text{``$2$ is even''}
\]

is a true statement, while

\[
\text{``$3$ is even''}
\]

is a false statement.

However,

\[
x+1
\]

is not a statement by itself, because it is not true or false. It is only an
expression. On the other hand,

\[
x+1>0
\]

can become a statement once we say what values of (x) we are considering.

In this chapter we introduce implication, necessary conditions, sufficient
conditions, logical equivalence, truth tables, and the contrapositive.

\subsection{Implication}

Let (A) and (B) be statements.

\begin{definition}\[Implication\]
The statement

\[
A \Rightarrow B
\]

means:

\[
\text{if $A$ is true, then $B$ must be true.}
\]

\end{definition}

We read this as:

\[
\text{``$A$ implies $B$''}
\]

or

\[
\text{``if $A$, then $B$.''}
\]

The implication $(A \Rightarrow B)$ does not say that (A) causes (B). It only
says that it is impossible for (A) to be true while (B) is false.

So the meaning of

\[
A \Rightarrow B
\]

is:

\[
\text{there is no case where $A$ is true and $B$ is false.}
\]

When (A) is false, the implication makes no demand on (B). For this reason,
in classical logic the implication $(A \Rightarrow B)$ is considered true
whenever (A) is false.

\subsection{Necessary and Sufficient Conditions}

Implication allows us to describe when one condition guarantees another.

\begin{definition}\[Sufficient Condition\]
A statement (A) is a \emph{sufficient condition} for (B) if

\[
A \Rightarrow B.
\]

This means that (A) being true is enough to guarantee that (B) is true.
\end{definition}

In other words, if we know (A), then we automatically know (B).

\begin{definition}\[Necessary Condition\]
A statement (B) is a \emph{necessary condition} for (A) if

\[
A \Rightarrow B.
\]

This means that whenever (A) is true, (B) must also be true. Without (B),
(A) cannot happen.
\end{definition}

Therefore, the same implication

\[
A \Rightarrow B
\]

has two interpretations:

\begin{itemize}
\item (A) is sufficient for (B);
\item (B) is necessary for (A).
\end{itemize}

This is one of the most important ideas in logic.

\subsection{Example: Divisibility}

Let

\[
A = \text{``$n$ is divisible by $6$''},
\]

and

\[
B = \text{``$n$ is divisible by $3$''}.
\]

If (n) is divisible by (6), then (n) is divisible by (3). Therefore,

\[
A \Rightarrow B.
\]

So:

\begin{itemize}
\item being divisible by (6) is sufficient for being divisible by (3);
\item being divisible by (3) is necessary for being divisible by (6).
\end{itemize}

However, being divisible by (3) is not sufficient for being divisible by (6).
For example,

\[
9
\]

is divisible by (3), but it is not divisible by (6).

So we have:

\[
A \Rightarrow B,
\]

but not

\[
B \Rightarrow A.
\]

\subsection{Logical Equivalence}

Sometimes two statements imply each other.

\begin{definition}\[Logical Equivalence\]
Statements (A) and (B) are \emph{logically equivalent} if

\[
A \Rightarrow B
\]

and

\[
B \Rightarrow A.
\]

In this case we write

\[
A \Leftrightarrow B.
\]

\end{definition}

We read this as:

\[
\text{``$A$ if and only if $B$.''}
\]

Logical equivalence means that (A) and (B) always have the same truth value.
Either both are true or both are false.

Equivalently,

\[
A \Leftrightarrow B
\]

means that (A) is both necessary and sufficient for (B), and (B) is both
necessary and sufficient for (A).

\subsection{Truth Tables}

Truth tables allow us to study logical statements by checking all possible truth
values.

\subsubsection{Implication}

The truth table for $(A \Rightarrow B)$ is:

\begin{center}
\begin{tabular}{c|c|c}
(A) & (B) & $(A \Rightarrow B)$ \\
\hline
T & T & T \\
T & F & F \\
F & T & T \\
F & F & T \\
\end{tabular}
\end{center}

The only situation where $(A \Rightarrow B)$ is false is when (A) is true and
(B) is false.

So the implication

\[
A \Rightarrow B
\]

is equivalent to saying:

\[
\text{not both $A$ true and $B$ false.}
\]

\subsubsection{Equivalence}

The truth table for $(A \Leftrightarrow B)$ is:

\begin{center}
\begin{tabular}{c|c|c}
(A) & (B) & $(A \Leftrightarrow B)$ \\
\hline
T & T & T \\
T & F & F \\
F & T & F \\
F & F & T \\
\end{tabular}
\end{center}

Equivalence is true exactly when (A) and (B) have the same truth value.

\subsection{Converse, Inverse, and Contrapositive}

Given an implication

\[
A \Rightarrow B,
\]

we can form three related statements.

\begin{definition}\[Converse\]
The \emph{converse} of $(A \Rightarrow B)$ is

\[
B \Rightarrow A.
\]

\end{definition}

\begin{definition}\[Inverse\]
The \emph{inverse} of $(A \Rightarrow B)$ is

\[
\neg A \Rightarrow \neg B.
\]

\end{definition}

\begin{definition}\[Contrapositive\]
The \emph{contrapositive} of $(A \Rightarrow B)$ is

\[
\neg B \Rightarrow \neg A.
\]

\end{definition}

The converse and inverse are not usually equivalent to the original implication.
However, the contrapositive is always equivalent to the original implication.

\begin{theorem}\[Contrapositive\]
The implication

\[
A \Rightarrow B
\]

is logically equivalent to its contrapositive

\[
\neg B \Rightarrow \neg A.
\]

\end{theorem}

\begin{proof}
The implication $(A \Rightarrow B)$ is false only in one case: when (A) is true
and (B) is false.

Now consider the contrapositive:

\[
\neg B \Rightarrow \neg A.
\]

This implication is false only when $(\neg B)$ is true and $(\neg A)$ is false.
But $(\neg B)$ being true means that (B) is false, and $(\neg A)$ being false
means that (A) is true.

So the contrapositive is false exactly when (A) is true and (B) is false.

Therefore, $(A \Rightarrow B)$ and $(\neg B \Rightarrow \neg A)$ are false in
exactly the same situation. Hence they have the same truth value in every case,
so they are logically equivalent.
\end{proof}

\subsection{Examples}

\subsubsection{Geometric Example}

Let

\[
A = \text{``It is a square''},
\]

and

\[
B = \text{``It is a rectangle with equal sides''}.
\]

Then:

\begin{itemize}
\item $(A \Rightarrow B)$: every square is a rectangle with equal sides;
\item $(B \Rightarrow A)$: every rectangle with equal sides is a square.
\end{itemize}

Therefore,

\[
A \Leftrightarrow B.
\]

So being a square is equivalent to being a rectangle with equal sides.

\subsubsection{Number Theory Example}

Let

\[
A = \text{``$n$ is divisible by $6$''},
\]

and

\[
B = \text{``$n$ is divisible by $3$''}.
\]

Then:

\[
A \Rightarrow B
\]

is true, because every number divisible by (6) is also divisible by (3).

However,

\[
B \Rightarrow A
\]

is false, because (9) is divisible by (3), but not by (6).

So:

\begin{itemize}
\item (A) is sufficient for (B);
\item (B) is necessary for (A);
\item (A) and (B) are not logically equivalent.
\end{itemize}

\subsubsection{Example with Inequalities}

Let

\[
A = \text{``$x>2$''},
\]

and

\[
B = \text{``$x>0$''}.
\]

Assume (x) is a real number. If (x>2), then certainly (x>0). Therefore,

\[
A \Rightarrow B.
\]

So (x>2) is sufficient for (x>0), and (x>0) is necessary for (x>2).

But the converse is not true. For example,

\[
x=1
\]

satisfies (x>0), but does not satisfy (x>2).

Thus,

\[
A \Rightarrow B,
\]

but not

\[
B \Rightarrow A.
\]

\subsection{The Importance of the Domain}

When working with logical statements in mathematics, we must know the domain of
the variables.

For example, the statement

\[
x^2 > 0
\]

depends on what values of (x) are allowed.

If (x) is a nonzero real number, then

\[
x^2>0
\]

is true.

But if (x=0) is allowed, then the statement is false for (x=0).

Therefore, before deciding whether an implication is true or false, we must know
the universe of objects we are discussing.

\subsection{Exercises}

\begin{enumerate}
\item Determine whether the following implication is true or false:

```
\[
\text{``If $n$ is prime, then $n$ is odd.''}
\]


If it is false, give a counterexample.

\item Determine whether the following implication is true or false:


\[
\text{``If $n$ is prime and $n>2$, then $n$ is odd.''}
\]


Identify the sufficient condition and the necessary condition.

\item Let


\[
A = \text{``$x>2$''},
\qquad
B = \text{``$x>0$''}.
\]


Write the converse, inverse, and contrapositive of \(A \Rightarrow B\).

\item Give an example of a condition that is necessary but not sufficient.

\item Give two conditions that are both necessary for a third condition, but
neither one alone is sufficient.

\item Construct the truth table for the statement


\[
(A \Rightarrow B) \land (\neg A \Rightarrow B).
\]


What simpler statement is it logically equivalent to?

\item Let


\[
A = \text{``$n$ is divisible by $4$''},
\qquad
B = \text{``$n$ is even''}.
\]


Decide whether \(A \Rightarrow B\) is true. Then decide whether
\(B \Rightarrow A\) is true.

\item Give an example of two statements \(A\) and \(B\) such that


\[
A \Leftrightarrow B.
\]
```

\end{enumerate}
